\documentclass[11pt,a4paper]{article}
\usepackage[T1]{fontenc}
\usepackage{lmodern,amsmath,amssymb,amsthm,microtype,booktabs}
\usepackage[margin=27mm]{geometry}
\usepackage[colorlinks=true,urlcolor=blue,linkcolor=blue]{hyperref}
\newtheorem{proposition}{Proposition}
\newcommand{\dd}{\mathop{}\!\mathrm d}
\title{Gram factors and transverse energy fluctuations\\in the recoil comparison}
\author{Alan Day \and Chengheng\thanks{AI research assistant (OpenAI).}}
\date{Working research note --- 30 October 2026}
\begin{document}\maketitle
\begin{abstract}
A symmetric energy slice is not the energy kernel of an antipodal
time correlator: two other energy differences must be integrated.
We factor the exact angle-Gram determinant in the CLP configuration,
then compute the transverse Hessian of the CSTY classical pure-recoil
action for equal probe masses. Conditionally using the geometric
one-loop factor, these two contributions supply the energy dependence
required by the previous recoil prefactor test. The full continued
Virasoro amplitude, its physical insertions and outer contour are not
computed. The analysis is an author self-reviewed draft, with no claim
of independent certification or priority.
\end{abstract}

\section{What the existing definitions provide}
Use de Sitter radius one, observer proper time \(t\), and
\(\kappa>0\). CLP~\cite{CLP}, (3.43)--(3.47), defines a fixed-topology
graph amplitude of the form
\begin{equation}
 Z_{\mathrm{tet}}
 =\left|\frac{\{6j\}_{b_+}}
 {\sqrt{C_{123}^+C_{156}^+C_{246}^+C_{345}^+}}\right|^2,
 \qquad b_+^2=i/\kappa. \label{eq:graph}
\end{equation}
The internal \(6j\) contour is specified by continuation. This is
not yet a complete identification with the canonically normalized,
fixed-average-mass mixed-representation correlator of CSTY~\cite{CSTY}.
In particular, cancellation of the order-\(\kappa\) vertex exponent
does not cancel the order-one vertex factors.

Following (3.64), CLP invokes a chiral one-loop Gram factor,
\[
 (\det G)^{-1/4},
\]
together with vertex corrections.
The cited Meng--Yang theorem~\cite{MY}, Theorem 1.3, is for fixed
hyperideal hyperbolic dihedral angles (with its stated conditions on
\(\arg b\)). Its \(O(b^2)\) remainder is not a uniform theorem for the
degenerating spherical configuration below. Our Gram identities are
exact algebra; using them as the whole continued \(6j\) one-loop
asymptotics remains conditional.

\section{An exact factorization and two different limits}
Use CLP's vertex triples \(123,156,246,345\). The angle Gram matrix
with face indices opposite these four vertices is
\begin{equation}
 G=\begin{pmatrix}
 1&-\cos\theta_4&-\cos\theta_5&-\cos\theta_6\\
 -\cos\theta_4&1&-\cos\theta_3&-\cos\theta_2\\
 -\cos\theta_5&-\cos\theta_3&1&-\cos\theta_1\\
 -\cos\theta_6&-\cos\theta_2&-\cos\theta_1&1
 \end{pmatrix}. \label{eq:gram}
\end{equation}
This labeling satisfies CLP Appendix B's discriminant identity
\(q_1^2-4q_0q_2=16\det G\).
Set
\[
 \theta_1=\theta_4=H+d,\quad \theta_3=\theta_6=H-d,\quad
 \theta_2=\pi-\beta',\quad\theta_5=\pi-\beta,
 \quad d=\frac{\pi\omega}{2\kappa},\quad
 (\beta,\beta')=\frac{\pi}{\kappa}(\nu,\nu').
\]
Here \(H\) is a cone half-angle, not the limiting special function
used in the previous note.

\begin{proposition}
Writing \(D(\omega)=\det G\), one has exactly, also for complex angles,
\begin{align}
 D(\omega)=16&
 \left[\cos^2\frac\beta2\cos^2\frac{\beta'}2-\cos^2H\cos^2d\right]
 \nonumber\\[-2pt]
 &\times
 \left[\sin^2\frac\beta2\sin^2\frac{\beta'}2-\sin^2H\sin^2d\right].
 \label{eq:factor}
\end{align}
\end{proposition}
\begin{proof}
Put \(a=\cos(H+d)\), \(b=\cos(H-d)\), \(u=-\cos\beta'\),
\(v=-\cos\beta\). Splitting the matrix into the two eigenspaces of
the exchange of its first and third, and second and fourth indices gives
\[
 D=\big[(1-u)(1-v)-(a+b)^2\big]
    \big[(1+u)(1+v)-(a-b)^2\big].
\]
The sum and difference formulas for \(a,b\) give \eqref{eq:factor}.
\end{proof}

At fixed \(H\) away from \(0,\pi\), let
\[
 A=\frac{\pi\nu\nu'}{2\kappa\sin H},\qquad \omega=A z.
\]
As \(\nu/\kappa,\nu'/\kappa\to0\), Taylor expansion of the two
factors gives
\begin{equation}
 \frac{D(Az)}{D(0)}=1-z^2+o(1), \label{eq:fixed}
\end{equation}
uniformly for bounded \(z\). Relative use excludes \(z=\pm1\).
For example, the unnormalized leading expression is
\(4\pi^2\sin^4H\,(A^2-\omega^2)/\kappa^2\).
This statement concerns the exact determinant, not the \(6j\) remainder.

A different limit takes \(H=\pi(1-M/\kappa)\) at fixed
\(M,\nu,\nu',\omega\). Direct expansion of \eqref{eq:factor} gives
\begin{equation}
 D(\omega)=\frac{\pi^6}{4\kappa^6}
 (4M^2-\nu^2-\nu'^2+\omega^2)
 (\nu^2\nu'^2-4M^2\omega^2)+O(\kappa^{-8}).
 \label{eq:recoil}
\end{equation}
If \(B=4M^2-\nu^2-\nu'^2>0\), \(A_0=\nu\nu'/(2M)\), then
\begin{equation}
 \lim_{\kappa\to\infty}\frac{D(\omega)}{D(0)}
 =\left(1+\frac{\omega^2}{B}\right)
  \left(1-\frac{\omega^2}{A_0^2}\right). \label{eq:recoilratio}
\end{equation}
Only in a further heavy-observer limit
\(\nu/M,\nu'/M\to0\), at bounded \(\omega/A_0\), may the first
factor be replaced by one. Thus \eqref{eq:fixed} is not simply a
uniform fixed-angle formula evaluated at \(H=\pi\).

If the chiral Gram factor is valid, its product with the other chirality,
on the real Euclidean germ, has normalized factor
\begin{equation}
 R_G(\omega)=\left(\frac{D(0)}{D(\omega)}\right)^{1/2}.
 \label{eq:RG}
\end{equation}
Its continuation uses that germ and the root equal to one at zero.
It is not defined by taking a fresh pointwise modulus at complex
\(\omega\). In the simplified limit,
\(R_G(iA_0 y)\to(1+y^2)^{-1/2}\). At the outer saddle
\(y=\sinh t\) this supplies only \((\cosh t)^{-1}\).
The previous total prefactor test required \((\cosh t)^{-2}\).

\section{The antipodal observable integrates two more energies}
CSTY (2.3) fixes the average of four observer segment masses.
Choose coordinates
\begin{equation}
 \begin{aligned}
 m_1&=\bar M+\omega/2+u,&m_3&=\bar M+\omega/2-u,\\
 m_2&=\bar M-\omega/2+v,&m_4&=\bar M-\omega/2-v.
 \end{aligned} \label{eq:coords}
\end{equation}
The absolute Jacobian is four, and the average is \(\bar M\).
For the antipodal times of CSTY (2.73), the four cyclic Euclidean
gaps are \(\pi/2-it,\pi/2+it,\pi/2-it,\pi/2+it\). Consequently
\begin{equation}
 \sum_i(\theta_i-\theta_{i-1})m_i=2\pi\bar M-2it\omega.
\end{equation}
Performing the average-mass delta integration in the defined transform
therefore gives
\begin{equation}
 \mathcal F(t\mid M)=\frac4{Z_{\mathrm{sphere}}(M)}
 \int\dd\omega\,e^{-2it\omega}
 \int_{\mathcal D_\omega}\dd u\,\dd v\,
 \boldsymbol{\mathcal F}(m_1,m_2,m_3,m_4).
 \label{eq:push}
\end{equation}
Domains and continued contours are inherited from the original
definition, not chosen anew. The thermal \(e^{-2\pi M}\) cancels
the constant phase above. Any Fourier convention constant is
independent of \(t\). CSTY footnote 19 explicitly notes the integration
over energy variables orthogonal to \(\omega\).

The restricted slice \(u=v=0\) cannot replace the inner integral.
For example the smooth local kernels
\(K(\omega)e^{-a(\omega)(u^2+v^2)/2}\), for different positive even
\(a\), have the same slice and different marginals \(2\pi K/a\).
This is an information-loss example, not a family of proposed physical
theories. At a nondegenerate local saddle of an action \(I\), the
missing factor is \(2\pi/\sqrt{\det I_{\perp}''}\), with its branch
and phase fixed by the transverse steepest-descent cycle.
This transverse integration is separate from the internal \(6j\)
Gaussian already included in a chiral one-loop determinant.

\section{Equal-probe transverse Hessian from the classical action}
Now specialize to \(\nu'=\nu>0\). This section uses the classical
opposite-side pure-recoil saddle of CSTY (2.54), (B.5)--(B.7):
\[
 I=\sum_{i<j}\mathbf m_{ij}d_{ij},\qquad
 \mathbf m=\begin{pmatrix}
 0&m_2&\nu&m_1\\m_2&0&m_3&\nu\\
 \nu&m_3&0&m_4\\m_1&\nu&m_4&0
 \end{pmatrix}.
\]
The angles \(d_{ij}\) follow from the dual flat tetrahedron, and are
continued from its nondegenerate real branch.
Assume \(B=4M^2-2\nu^2>0\) and \(|\omega|<A_0=\nu^2/(2M)\).
The transverse saddle is \(u=v=0\).

\begin{proposition}
On this Euclidean branch,
\begin{equation}
 I_{\perp}''=h(\omega)\mathbf1_2,\qquad
 h(\omega)=
 \frac{4\sqrt{(\nu^4-4M^2\omega^2)(4M^2-2\nu^2+\omega^2)}}
 {(4M^2-\nu^2)(\nu^2-\omega^2)}.
 \label{eq:h}
\end{equation}
The root is positive at zero. The formula extends analytically until
a root, pole or branch obstruction is reached.
\end{proposition}
\begin{proof}
Stationarity of the vertex positions implies
\(\dd I=\sum d_{ij}\dd\mathbf m_{ij}\), equivalently the flat-dual
Schl\"afli identity cancels the mass-weighted angle variations.
Thus \(I_u=d_{14}-d_{23}\), \(I_v=d_{12}-d_{34}\), both zero
on the symmetric slice. Let
\(D_*=(4M^2-\nu^2)(\nu^2-\omega^2)\).
Differentiating CSTY's explicit formula for \(C=\cos d_{14}\) gives
\begin{align}
 C&=\frac{8M^3\omega+4M^2\omega^2-4M\nu^2\omega
              +2M\omega^3-\nu^4}{D_*},\\
 \left.\partial_u C\right|_0
 &=\frac{2(2M+\omega)(4M^2\omega^2-\nu^4)
                   (4M^2-2\nu^2+\omega^2)}{D_*^2},
 \qquad \left.\partial_v C\right|_0=0,\\
 1-C^2&=\frac{(2M+\omega)^2(\nu^4-4M^2\omega^2)
                   (4M^2-2\nu^2+\omega^2)}{D_*^2}.
\end{align}
Exchange of the opposite masses reverses the \(u\) derivative of
\(d_{23}\). Hence \(I_{uu}=-2\partial_u C/\sqrt{1-C^2}\),
giving \eqref{eq:h}. The other diagonal entry follows with
\(\omega\to-\omega\), and the mixed derivative vanishes.
\end{proof}

The normalized transverse Gaussian factor is therefore
\begin{equation}
 R_\perp(\omega)=\frac{h(0)}{h(\omega)}
 =\frac{1-\omega^2/\nu^2}
 {\sqrt{1-\omega^2/A_0^2}\sqrt{1+\omega^2/B}}.
 \label{eq:RT}
\end{equation}
Combining it with the pure-recoil Gram ratio in
\eqref{eq:recoilratio} yields
\begin{equation}
 R_G R_\perp=
 \frac{1-\omega^2/\nu^2}
 {(1-\omega^2/A_0^2)(1+\omega^2/B)}.
 \label{eq:combined}
\end{equation}
For \(\nu/M\to0\), at \(\omega=iA_0 y\) with \(y\) in a fixed
bounded interval, this tends uniformly to \((1+y^2)^{-1}\).
This is precisely the energy dependence required by the previous
prefactor diagnostic. It comes from geometry and marginalization,
not fitting a function to a target correlator.

This is a \emph{conditional reconciliation of two local one-loop
pieces}. We have not proved that these pieces exhaust the physical
amplitude. Vertex factors, physical insertions, the other saddle and
the outer contour remain. In particular, positivity of the Euclidean
local Hessian is not a proof that its continued cycle contributes to the
desired Lorentzian observable.

The transverse Gaussian width is of order \(\sqrt M\), since
\(h(0)\sim2/M\). Resolving it well within the probe scale suggests
\(\sqrt M\ll\nu\ll M\). Together with the previous reduced-function
window \(A_0\cosh t\ll\nu^{2/3}\), for bounded \(t\), a compatible
power hierarchy is \(\nu^{4/3}\ll M\ll\nu^2\), with large \(\nu\).
This is a compatibility check of necessary scales, not a uniform
remainder theorem for the full integral or its finite-\(\kappa\)
degeneration.

\section{What remains unspecified or uncomputed}
The exact \(C_{ijk}\) functions in \eqref{eq:graph} are available;
their joint-limit contribution is a calculation still to do.
The physical two-point insertion normalization in CLP (3.37),
however, is prescribed only to leading order at fixed probe data.
It does not alone determine a joint growing-probe limit. For example,
multiplying that leading prescription by
\[
 1+c\,\frac{\nu^2}{\kappa+\nu^2}\,
             \frac{\omega^2}{1+\omega^2},\qquad c>0,
\]
preserves its fixed-\(\nu,\omega\) relative \(O(\kappa^{-1})\)
accuracy and real-label exchange symmetry, but can change a joint
limit by order one. This demonstrates insufficiency of that asymptotic
input alone; no full consistency of such an arbitrary completion is
claimed.

Likewise, the specified internal \(6j\) contour is not the complete
external three-energy integration cycle. CLP footnote 12 restricts
the \(i\mathbb R\) prescription to the reduced action.
A continued Euclidean modulus square also cannot be evaluated as a
new modulus square of complex energies; the two chiral analytic factors
must be continued with their original branches.

The next calculable step is the vertex/leg contribution in the same
equal-probe hierarchy, retaining the two transverse energies before
reduction. It must be accompanied by a justified physical insertion
normalization and an outer cycle. If these definitions are not fixed,
the output must remain a conditional graph calculation rather than a
physical finite-\(G\) OTOC.

\section{Evidence and reproducibility}
The script \texttt{code/gram\_prefactor.py} saves 83 checks:
symbolic Gram factorization, the mass Jacobian and phase;
five discriminant checks (including complex angles);
nine growing-probe fixed-angle cases; six pure-recoil cases;
six Gram-weighted inverse integrals; and nine transverse Hessians.
The last are differentiated directly from the six-angle action and
compared with \eqref{eq:h}. Working precision is 65 decimal digits;
the indicated determinant ratios and inverse integrals are repeated
at 95 and 90 digits. This is arithmetic evidence, not interval
certification, a proof of a full \(6j\) asymptotic, or a finite-\(G\)
calculation. Existing frozen calculations were not altered.

\begin{thebibliography}{9}
\bibitem{CLP} J. Chandra, D. Liska and B. Post,
\emph{3D de Sitter from complex Virasoro TQFT: Observers and knots},
\href{https://arxiv.org/html/2610.30247v1}{arXiv:2610.30247v1},
Sections 3.2--3.3, footnote 12, Appendices B--C.
\bibitem{CSTY} Y. Chen, D. Stanford, H. Tang and Z. Yang,
\emph{Negative shocks versus static patch holography},
\href{https://arxiv.org/html/2607.14042v1}{arXiv:2607.14042v1},
(2.3), (2.54), Section 2.6, footnote 19 and Appendix B.1.
\bibitem{MY} Y. Meng and T. Yang,
\emph{Asymptotics of complex \(b\)-\(6j\) symbols},
\href{https://arxiv.org/html/2606.03653v1}{arXiv:2606.03653v1},
Theorem 1.3. Read online 30 October 2026; not locally cached.
\end{thebibliography}
\end{document}
